% =========================================================
% 4. unlink() ORM Method
% =========================================================
\section{\texttt{unlink()} ORM Method}

\begin{frame}[standout]
  \texttt{unlink()} ORM Method
\end{frame}

\begin{frame}[fragile]{\texttt{UNLINK()} METHOD}
  \begin{itemize}
    \item The \texttt{unlink()} method is an ORM method used to \textbf{delete} one or more
          existing records from an Odoo model.
    \item In database terms, \texttt{unlink()} is similar to SQL \texttt{DELETE}.
    \item The method returns a boolean value: \texttt{True} when deletion succeeds.
  \end{itemize}
  \begin{block}{Method Syntax}
    The standard syntax of the \texttt{unlink()} method is:
\begin{lstlisting}
def unlink(self):
    return super().unlink()
\end{lstlisting}
    This method has two important parts:
    \begin{enumerate}
      \item \texttt{unlink}
      \item \texttt{self}
    \end{enumerate}
  \end{block}
\end{frame}

\begin{frame}[fragile]{BREAKING DOWN THE METHOD SIGNATURE}
  \begin{itemize}
    \item \texttt{unlink}: The ORM method responsible for deleting records.
    \item \texttt{self}: A recordset of one or more records that will be deleted.
  \end{itemize}
\begin{lstlisting}
student = self.env['student.student'].browse(15)
print(student)
\end{lstlisting}
  Output is:
\begin{lstlisting}
student.student(15,)
\end{lstlisting}
  \begin{itemize}
    \item \texttt{unlink()} does \textbf{not} take a \texttt{vals} dictionary.
    \item It only needs the recordset that should be removed.
  \end{itemize}
\end{frame}

\begin{frame}[fragile]{WHAT DOES \texttt{UNLINK()} RETURN?}
  \begin{itemize}
    \item The \texttt{unlink()} method returns \texttt{True} on success.
  \end{itemize}
\begin{lstlisting}
student = self.env['student.student'].browse(15)
result = student.unlink()
print(result)
\end{lstlisting}
  Output is:
\begin{lstlisting}
True
\end{lstlisting}
  \begin{itemize}
    \item After \texttt{unlink()}, the record no longer exists in the database.
    \item Related one2many / many2many cleanup is handled by the ORM according to
          the field \texttt{ondelete} behavior.
  \end{itemize}
\end{frame}

\begin{frame}[fragile]{DELETING ONE RECORD VS MANY RECORDS}
  \begin{enumerate}
    \item \textbf{One record:}
\begin{lstlisting}
student = self.env['student.student'].browse(15)
student.unlink()
\end{lstlisting}
    \item \textbf{Many records:}
\begin{lstlisting}
students = self.env['student.student'].search([
    ('active', '=', False),
])
students.unlink()
\end{lstlisting}
  \end{enumerate}
  \begin{itemize}
    \item Prefer one \texttt{unlink()} on a recordset over deleting records in a Python loop.
    \item Access rights and record rules still apply before deletion.
  \end{itemize}
\end{frame}

\begin{frame}[fragile]{METHOD CALL AND OVERRIDE PATTERN}
  \textbf{Method Call}
\begin{lstlisting}
student.unlink()
\end{lstlisting}
  \begin{itemize}
    \item Here, \texttt{.unlink()} is the method call.
  \end{itemize}
  \textbf{Common Override Pattern}
\begin{lstlisting}
def unlink(self):
    for record in self:
        if record.state == 'done':
            raise UserError("Done students cannot be deleted.")
    return super().unlink()
\end{lstlisting}
  \begin{itemize}
    \item Overrides are often used to block unsafe deletions.
  \end{itemize}
\end{frame}

\begin{frame}[fragile]{IMPORTANT NOTES ABOUT \texttt{UNLINK()}}
  \begin{itemize}
    \item Deletion is permanent for normal models (unless you use archive/\texttt{active}).
    \item Many business models prefer \textbf{archiving} instead of deleting:
  \end{itemize}
\begin{lstlisting}
student.write({'active': False})  # soft delete / archive
\end{lstlisting}
  \begin{itemize}
    \item Foreign-key constraints may raise errors if other records still depend on
          the record you are deleting.
    \item Always check business rules before calling \texttt{unlink()} in production code.
  \end{itemize}
\end{frame}
